\documentclass[11pt]{article}
\usepackage{mathblatt}
% Probe Serie 08.10.2026, Paket Pythagoras: Lösungen zu Lernblatt 1 (Kennung QH4), von Hand gesetzt.
% Je Teilaufgabe eine Zeile: links fett das Ergebnis mit Einheit, rechts klein Ansatz ⇒ Wert; exakt vor gerundet.
% Alle Werte mit sympy nachgerechnet.
\newcommand{\kennung}{QH4}
\begin{document}
\pfheftstil
\fancyfoot[C]{{\footnotesize\color{mbgrau}{\scriptsize \kennung}\hfill\thepage}}
\pfheftkopf{Die lange Seite \textperiodcentered{} Lösungen}{Lernblatt 1}

\begin{pfloesung}{}
\lz{1a)}{die längste Seite}{Hypotenuse liegt dem rechten Winkel gegenüber}{}
\lz{b)}{9 und 16 Kästchen}{$3^2 = 9$, \ $4^2 = 16$}{}
\lz{c)}{25 Kästchen}{$5^2 = 25$}{}
\lz{d)}{die Quadrate}{$9 + 16 = 25$; \ die Seiten nicht: $3 + 4 = 7 \ne 5$}{}
\end{pfloesung}
\begin{pfloesung}{}
\lz{2a)}{$r$}{gegenüber dem rechten Winkel}{}
\lz{b)}{$m$}{gegenüber dem rechten Winkel}{}
\lz{c)}{$e$}{gegenüber dem rechten Winkel}{}
\end{pfloesung}
\begin{pfloesung}{}
\lz{3.}{die dritte Aussage}{Kathetenquadrate zusammen = Hypotenusenquadrat}{}
\end{pfloesung}
\begin{pfloesung}{}
\lz{4a)}{$u^2 + w^2 = v^2$}{Hypotenuse $v$ allein}{}
\lz{b)}{$d^2 + h^2 = s^2$}{Hypotenuse $s$ allein}{}
\end{pfloesung}
\begin{pfloesung}{}
\lz{5.}{$z = \sqrt{x^2 + y^2}$}{$z^2 = x^2 + y^2$ \ $\Rightarrow$ \ Wurzel ziehen}{}
\end{pfloesung}
\begin{pfloesung}{}
\lz{6a)}{die Hypotenuse}{beide Katheten gegeben}{}
\lz{b)}{36 und 64}{$6^2 = 6 \cdot 6 = 36$, \ $8^2 = 64$}{}
\lz{c)}{$c^2 = 36 + 64$}{$c^2 = 100$}{}
\lz{d)}{$c = 10$ cm}{$c = \sqrt{100}$ \ $\Rightarrow$ \ $10$ cm}{}
\end{pfloesung}
\begin{pfloesung}{}
\lz{7.}{$c = 39$ cm}{$c^2 = 36^2 + 15^2 = 1\,296 + 225 = 1\,521$ \ $\Rightarrow$ \ $\sqrt{1\,521} = 39$}{}
\end{pfloesung}
\begin{pfloesung}{}
\lz{8.}{$c = \sqrt{97} \approx 9{,}8$ cm}{$c^2 = 4^2 + 9^2 = 16 + 81 = 97$ \ $\Rightarrow$ \ $\sqrt{97} \approx 9{,}85$}{}
\end{pfloesung}
\begin{pfloesung}{}
\lz{9.}{$d = 51$ cm}{Hypotenuse: die Diagonale; \ $d^2 = 45^2 + 24^2 = 2\,601$ \ $\Rightarrow$ \ $51$}{}
\end{pfloesung}
\begin{pfloesung}{}
\lz{10.}{$x = 2\sqrt{7\,289} \approx 170{,}8$ cm}{gesucht: die schräge Rampe; \ $x^2 = 170^2 + 16^2 = 29\,156$ \ $\Rightarrow$ \ $\sqrt{29\,156} \approx 170{,}75$}{}
\end{pfloesung}
\begin{pfloesung}{}
\lz{11.}{$c = 65$ cm}{Wurzel vergessen: $c = \sqrt{4\,225}$ \ $\Rightarrow$ \ $65$}{}
\end{pfloesung}
\begin{pfloesung}{}
\lz{12.}{$\overline{AB} = 10\sqrt{13} \approx 36{,}1$ m}{$\overline{AB}^2 = 12^2 + 34^2 = 144 + 1\,156 = 1\,300$ \ $\Rightarrow$ \ $\sqrt{1\,300} \approx 36{,}06$}{}
\end{pfloesung}
\end{document}
